\subsection{相对膜结构}\label{sec:17i}
设 $\pi:E\to S$ 为 Lefschetz 纤维化；我们先讨论 $S$ 紧的情形，稍后再处理带条带状端的版本。假设 $E$ 带有相对二次复体积形式 $\eta^2_{E/S}$，以 $\alpha_{E/S}$ 表示相对平方相位映射。再假设 $F$ 为 $E$ 上的 Lagrangian 边界条件。
\begin{definition*}
$F$ 的\emph{相对分次}是映射 $x\mapsto\alpha_{E/S}(TF_x^v)$ 的一个提升 $\alpha^\#:F\to\R$。$F$ 上的\emph{相对膜结构}是二元组 $(\alpha^\#,P^\#)$，由相对分次与相对 Pin 结构组成；后者就是向量丛
\SourcePage{267}
$TF^v\to F$ 上的一个 Pin 结构 $P^\#$。
\end{definition*}
这类结构存在的障碍可作如下理解：在每个边界圆周 $C$ 上选一点 $z$。每个 $F_z\subset E_z$ 都必须能关于 $\eta^2_{E_z}$ 赋予分次；此外，若 $h_C$ 为沿该边界的单值变换，并如\secref{15c}中那样带有自然分次，则其像 $h_C(F_z)$ 必须与 $F_z$ 作为分次 Lagrangian 子流形相等。类似地，每个 $F_z$ 必须有 Pin 结构 $P_z^\#$，使 $(h_C|F_z)^*P_z^\#$ 同构于 $P_z^\#$；选定一个这样的同构，就确定了 $F|C$ 的相对 Pin 结构。

\StatementNumber{14}\begin{example}\label{st:17.14}
取\stref{例题}{17.3}中的局部模型，为简单起见令 $r=1$。标准相对二次复体积形式是 $\eta_E=dx_1\wedge\cdots\wedge dx_{n+1}$ 与 $\eta_S=dz$ 的商的平方。对于每个 $x\in F_z$，有
\[
\alpha_{E/S}(TF_x^v)=\frac{\eta_E(v_1,\ldots,v_{n+1})^2}{\eta_S(iz\partial_z)^2},
\]
其中 $\{v_1,\ldots,v_{n+1}\}$ 为整个 $F$ 的切空间的一组标准正交基。具体地，可选 $v_1,\ldots,v_n$ 为 $\sqrt zS^n\subset\sqrt z\R^{n+1}$ 的切向量，并取 $v_{n+1}=ix\in i\sqrt z\R^{n+1}$，从而 $\alpha_{E/S}(TF_x^v)=z^{n-1}$。因此，$n=1$ 时 $F$ 存在相对分次，而 $n>1$ 时不存在。\znote{259}{$\alpha_{E/S}$ 取值于 $S^1$，相对分次是它的实值提升，绕临界值一圈的卷绕数就是障碍。此处 $\alpha_{E/S}(TF_x^v)=z^{n-1}$ 绕 $z=0$ 卷绕 $n-1$ 次，故 $n>1$ 时无整体提升；$n=1$ 时 $\alpha_{E/S}\equiv1$，常值提升即为分次。又因 $\alpha$ 是平方相位，其卷绕数是相对第一 Chern 数的两倍，这正是\stref{笔记}{17.12}中 $2c_1(M,V)=0$ 的来源。[来源：Seidel, Fukaya categories and Picard–Lefschetz theory, §17i。]}这反映了如下事实：任意 Dehn 扭转 $\tau_V$ 赋予其自然分次后，都会将中心球面 $V$ 的分次平移 $1-n$，见\bcite{125}引理 5.7。相对 Pin 结构则始终存在；事实上，$n=1$ 时有四种不同选择，$n>1$ 时有两种。
\end{example}

分次的存在保证：对任意边界值落在 $F$ 中的截面 $u$，丛对 $(u^*(TE^v),u^*(TF^v))$ 的相对第一 Chern 类为零。因此，假定正则性成立，模空间的维数就是 $n\chi(S)$。与前面一样，在 $\partial S$ 的每个分支上选一点，记为 $z_j$。给定相对 Pin 结构后，就得到 $\M_{E/S}$ 相对于各 $z_j$ 处求值映射之积的一个优选定向。这是\stref{命题}{11.13}的推论，只须将抽象膜结构的语言转换成当前的语言。于是，这些带优选相对定向的求值映射，通过 Poincaré 对偶表示一个上同调类
\begin{equation}\label{eq:17.22}
\Phi_{E/S,z_1,\ldots,z_k}\in H^{n(k-\chi(S))}\Bigl(\prod_jF_{z_j}\Bigr).
\end{equation}
取任意实线丛 $\beta\to\partial S$。可以通过 $P^\#\mapsto P^\#\otimes(\pi|F)^*\beta$ 改变相对 Pin 结构；这对\eref{17.22}的影响是将不变量乘以 $(-1)^i$，其中 $i$ 是 $\beta$ 非平凡的分支数。最后，这一整套讨论都可毫无困难地推广到移动边界条件。

对我们而言，\eref{17.22}最重要的特殊情形是 $S$ 为圆盘、纤维维数为 $2n=2$。后一个假设意味着，线性化算子 $D_{E/S,u}$ 是线丛上的 $\bar\partial$ 算子，因此\secref{11d}中的论证
\SourcePage{268}
适用。特别地，由\eref{11.13}可知，对每个 $u\in\M_{E/S}$，$\ker(D_{E/S,u})$ 都是一维的，并由 $u^*(TE^v)$ 的一个处处非零截面 $X$ 生成。在 $X$ 给出的平凡化下，$D_{E/S,u}$ 就是平凡线丛上具有常值边界条件的标准 $\bar\partial$ 算子。回顾\stref{引理}{11.12}的证明可见，若 $u^*(TF^v)$ 上的 Pin 结构平凡，则相对定向 $\ker(D_{E/S,u})\cong(TF_z)_{u(z)}$ 由映射 $X\mapsto X(z)$ 给出；否则，定向与之相反。这意味着，可按如下方式计算局部常值函数 $\Phi_{E/S,z}\in H^0(F_z)$：选一点 $x\in F_z$，然后计数满足 $u(z)=x$ 的截面 $u\in\M_{E/S}$。若 $u^*(TF^v)$ 上的 Pin 结构平凡，则赋予重数 $+1$；否则赋予 $-1$。

\StatementNumber{15}\begin{example}\label{st:17.15}
回到\stref{例题}{17.3}，假设 $n=1$。边界条件可显式参数化为
\[
S^1\times S^1\longrightarrow F,\qquad
(t_1,t_2)\longmapsto\bigl(e^{it_1/2}\cos(t_1/2+t_2),e^{it_1/2}\sin(t_1/2+t_2)\bigr).
\]
截面模空间是两个圆周 $\M_{E/S,\pm}$ 的不交并，它们由\eref{17.8}中取 $a=\tfrac12(e^{it},\pm ie^{it})$ 得到。在这些坐标下，沿边界的求值映射成为
\[
\begin{aligned}
\partial S\times\M_{E/S,+}&\longrightarrow F,&(s,t)&\longmapsto(t_1,t_2)=(s,t-s),\\
\partial S\times\M_{E/S,-}&\longrightarrow F,&(s,t)&\longmapsto(t_1,t_2)=(s,t).
\end{aligned}
\]
若固定 $F$ 中一点，比如 $(t_1,t_2)=(0,0)$，则每族中都恰有一个截面 $u_\pm$ 经过该点，即 $t=0$ 所对应的截面。注意，在同调中，圆周 $u_\pm(\partial S)\subset F$ 相差一份纤维 $F_z$。根据前面的一般性观察，可得如下结论：若选择的相对 Pin 结构在 $F_z$ 上\emph{非平凡}，则 $u_+$ 与 $u_-$ 的贡献相互抵消，不变量为零。否则，$\Phi_{E/S,z}=\pm2$；两个符号都会出现，对应的两个 Pin 结构相差 $\partial S$ 上非平凡线丛的拉回。\znote{260}{$u_+(\partial S)$ 与 $u_-(\partial S)$ 在 $F$ 中只相差一份纤维类 $F_z$，故相对 Pin 结构限制到 $u_\pm^*(TF^v)$ 上时恰相差一次符号；结合前面\stref{引理}{11.12}的观察（Pin 结构平凡给 $+1$、非平凡给 $-1$），两条截面的贡献相消。只有当 Pin 结构在 $F_z$ 上平凡时二者同号，总数才得 $\pm2$。[来源：Seidel, Fukaya categories and Picard–Lefschetz theory, §17i 与\stref{引理}{11.12}。]}
\end{example}

固定带角正合辛流形 $M$ 及其上的二次复体积形式 $\eta_M^2$。取一个带条带状端的 Lefschetz 纤维化，使每个端上的平凡化都以 $Z^\pm\times M$ 为模型。为它赋予一个可能移动的 Lagrangian 边界条件 $F$，以 $(L_{\zeta,0},L_{\zeta,1})$ 表示与每个端 $\zeta$ 相关的 Lagrangian 子流形对。现在假设，它们带有膜结构 $(L_{\zeta,0}^\#,L_{\zeta,1}^\#)$，且 $E$ 带有相对二次复体积形式 $\eta^2_{E/S}$，在每个端上都与 $\eta_M^2$ 一致。
\begin{definition*}
此时，$F$ 上的\emph{相容相对膜结构}由下列数据组成：相对分次 $\alpha^\#$，在每个端上等于各 $L_{\zeta,k}$ 的分次；相对 Pin 结构 $P^\#$；以及该结构与各 $L_{\zeta,k}$ 在端上所带的 Pin 结构之间的指定同构。
\end{definition*}
有了这种结构，便可对行列式线丛作与\secref{12c}类似的分析。当每个边界圆周恰含一个入向无穷远点时，结论最容易陈述；我们所有的应用都属于这种情形。此时，对
\SourcePage{269}
$\M_{E/S}(\{y_\zeta\})$ 中孤立点作带符号计数，就得到\eref{17.10}的精细版本，即上链映射
\[
C\Phi_{E/S}:\bigotimes_{\zeta_+\in\Sigma^+}CF^*(L_{\zeta_+,0}^\#,L_{\zeta_+,1}^\#)
\longrightarrow\bigotimes_{\zeta_-\in\Sigma^-}CF^*(L_{\zeta_-,0}^\#,L_{\zeta_-,1}^\#)[n(\chi(S)-|\Sigma^-|)],
\]
其中 $CF^*$ 是以任意域 $\K$ 为系数的 $\Z$ 分次 Floer 上链复形。还可以允许更一般的曲面，但此时要么必须为某些 $F_z$ 选定定向，见\eref{12.5}；要么必须像\eref{12.7}中那样，改变与端相关的某些膜结构。
